\section{Mathematical groundwork}
\label{sec:math}

This section fixes notation and collects the mathematics the rest of the survey
uses. We have written it to be self-contained for a reader who knows
undergraduate algebra and probability but not cryptography, because the
comparison this paper makes is between paradigms whose costs are determined by
quite different mathematical objects---a carry chain in one case, the growth of a
noise term in another---and a reader cannot weigh our measurements without both.
Where a statement is standard we give it with a proof anyway when the proof is
short and the constant matters to us later; where the proof is long or the result
is purely cited, we say so and give the reference.

\subsection{Notation}
\label{sec:math:notation}

\begin{notation}
\label{not:basic}
We write $\N$, $\Z$, $\Q$, $\Rr$ and $\Cc$ for the naturals, integers,
rationals, reals and complex numbers. For $m \in \N$ we write $\Z_m$ for the
ring of integers modulo $m$; we use $\Z_{2^k}$ for the machine-word rings and
$\Z_p$ or $\mathbb{F}_p$ for prime fields, writing the latter as a field when we
rely on invertibility. For a finite set $S$ we write $s \gets S$ for a draw from
the uniform distribution on $S$, and $s \gets D$ for a draw from a distribution
$D$. Vectors are lower-case bold, $\mathbf{a}$, and matrices upper-case bold,
$\mathbf{A}$. For $x \in \Rr$ we write $\floorr{x}$, $\ceil{x}$ and $\rnd{x}$
for rounding down, up, and to the nearest integer with ties resolved upward.
For $a \in \Z$ and odd or even modulus $m$ we write $[a]_m$ for the
representative of $a + m\Z$ in the \ib{centred} interval
$[-m/2, m/2)$, and $a \bmod m$ for the representative in $[0,m)$; the
distinction matters in Section~\ref{sec:math:fixed}, because a fixed-point value
is a centred representative while a machine word is the other one.
The security parameter is $\lambda$, the statistical parameter is $\kappa$, and
$\log$ is base two unless written otherwise. We write $\abs{S}$ for the
cardinality of a set and $\norm{\cdot}_\infty$ for the maximum-absolute-value
norm on a vector or on the coefficient vector of a polynomial. Parties are
$P_0,\dots,P_{n-1}$ and indices into parties are always modulo $n$, so that
$P_{i+1}$ makes sense for $i = n-1$.
\end{notation}

We reserve $\shr{x}$ for a secret sharing of $x$ in whichever scheme is in force,
and $\mathsf{ct}$ for a ciphertext. Overloading $\shr{\cdot}$ is deliberate: one
of the points of this survey is that a program can be written once against an
abstract sharing and executed under several concrete ones, and the notation
should not pretend otherwise. Where the concrete scheme matters we subscript it.

\subsection{Asymptotics, negligible functions and the two kinds of security}
\label{sec:math:negl}

Cryptographic statements are asymptotic: they say that an adversary's advantage
shrinks faster than any inverse polynomial in a parameter the honest parties may
increase. Making this precise is what lets us distinguish claims that are
\emph{unconditional} from claims that rest on a computational assumption, and the
distinction runs right through our comparison, because secret sharing with an
honest majority is information-theoretic while homomorphic encryption is not.

\begin{definition}
\label{def:negl}
A function $\varepsilon : \N \to \Rr_{\ge 0}$ is \ib{negligible} when for every
$c \in \N$ there is an $n_0$ such that $\varepsilon(n) < n^{-c}$ for all
$n > n_0$. We write $\negl(n)$ for an unspecified negligible function and
$\poly(n)$ for an unspecified positive polynomial. A function is
\ib{overwhelming} when it is $1 - \negl(n)$.
\end{definition}

\begin{proposition}
\label{prop:negl-closure}
Negligible functions are closed under addition and under multiplication by a
polynomial: if $\varepsilon_1,\varepsilon_2$ are negligible and $p$ is a
polynomial, then $\varepsilon_1 + \varepsilon_2$ and $p \cdot \varepsilon_1$ are
negligible.
\end{proposition}

\begin{proof}
Fix $c$. Choose $n_1$ with $\varepsilon_1(n) < n^{-(c+1)}$ for $n > n_1$ and
$n_2$ likewise for $\varepsilon_2$. For $n > \max(n_1,n_2,2)$ we have
$\varepsilon_1(n)+\varepsilon_2(n) < 2n^{-(c+1)} \le n^{-c}$, which proves the
first claim. For the second, let $p$ have degree $d$ and leading behaviour
bounded by $p(n) \le n^{d+1}$ for large $n$; choose $n_3$ with
$\varepsilon_1(n) < n^{-(c+d+1)}$ for $n > n_3$, so that
$p(n)\varepsilon_1(n) < n^{d+1}n^{-(c+d+1)} = n^{-c}$.
\end{proof}

Proposition~\ref{prop:negl-closure} is the reason a protocol may invoke a
subprotocol polynomially many times without losing security, and it is used
implicitly in every hybrid argument below.

\begin{definition}
\label{def:sd}
The \ib{statistical distance} between distributions $D_0,D_1$ over a countable
set $S$ is
\[
  \SD(D_0,D_1) \;=\; \tfrac{1}{2}\sum_{s \in S}\abs{\Pr[D_0 = s] - \Pr[D_1 = s]}
  \;=\; \max_{T \subseteq S}\;\bigl(\Pr[D_0 \in T] - \Pr[D_1 \in T]\bigr).
\]
Two families of distributions are \ib{statistically indistinguishable} when their
statistical distance is negligible in $\lambda$, and \ib{perfectly
indistinguishable} when it is zero.
\end{definition}

The two expressions in Definition~\ref{def:sd} agree because the maximising $T$
is $\{s : \Pr[D_0=s] > \Pr[D_1=s]\}$, and the positive and negative parts of the
difference each sum to the same value since both distributions have total mass
one. The second form is the operational one: it says that statistical distance is
exactly the advantage of the best possible distinguisher, \emph{including
computationally unbounded ones}. This is why a statistical guarantee needs no
assumption, and why we quote $\kappa$ separately from $\lambda$ throughout.

\begin{definition}
\label{def:ci}
Distribution ensembles $\{D_0(\lambda)\}$ and $\{D_1(\lambda)\}$ are
\ib{computationally indistinguishable}, written $D_0 \approx_c D_1$, when for
every probabilistic polynomial-time distinguisher $\mathcal{D}$,
\[
  \abs{\Pr[\mathcal{D}(1^\lambda, D_0(\lambda)) = 1] -
       \Pr[\mathcal{D}(1^\lambda, D_1(\lambda)) = 1]} \;=\; \negl(\lambda).
\]
\end{definition}

Statistical indistinguishability implies computational indistinguishability, and
the converse fails; the gap is where all of lattice cryptography lives.

\begin{lemma}[Hybrid argument]
\label{lem:hybrid}
Let $D_0,D_1,\dots,D_m$ be distribution ensembles with $m = \poly(\lambda)$, and
suppose $D_{j-1} \approx_c D_j$ for every $j$. Then $D_0 \approx_c D_m$.
\end{lemma}

\begin{proof}
Let $\mathcal{D}$ be a distinguisher and write
$p_j = \Pr[\mathcal{D}(D_j)=1]$. By the triangle inequality,
$\abs{p_0 - p_m} \le \sum_{j=1}^{m}\abs{p_{j-1}-p_j}$. Each term is negligible by
hypothesis, and there are polynomially many of them, so the sum is negligible by
Proposition~\ref{prop:negl-closure}.
\end{proof}

Lemma~\ref{lem:hybrid} is the workhorse of every simulation proof in this survey:
a simulator is exhibited by replacing one real message at a time with a simulated
one and arguing that each replacement is undetectable.

We will need one further probabilistic tool, for the arguments in which a mask
drawn from a wide interval hides a value drawn from a narrow one. It complements
Theorem~\ref{thm:mask}, which treats masking over a whole ring; the lemma below
treats masking by an interval, which is the situation in threshold decryption of a
lattice ciphertext (Section~\ref{sec:prelim:fhe}) and in the masked openings of
Section~\ref{sec:math:bitdec}.

\begin{lemma}[Smudging]
\label{lem:smudge}
Let $B, B' \in \N$ with $B' \ge 2^{\kappa} B$, let $e$ be any fixed integer with
$\abs{e} \le B$, and let $E \gets [-B',B']$ be uniform. Then
\[
  \SD\bigl(E,\; E + e\bigr) \;\le\; \frac{\abs{e}}{2B'+1} \;\le\; 2^{-\kappa}.
\]
Consequently, for any two integers $e_0,e_1$ of magnitude at most $B$, the
distributions $E+e_0$ and $E+e_1$ are at statistical distance at most
$2^{-\kappa+1}$.
\end{lemma}

\begin{proof}
Both $E$ and $E+e$ are uniform on integer intervals of the same cardinality
$2B'+1$, namely $[-B',B']$ and $[-B'+e, B'+e]$. Two uniform distributions on
finite sets $S_0, S_1$ of equal cardinality $N$ have statistical distance
$\abs{S_0 \setminus S_1}/N$: each point of $S_0\setminus S_1$ contributes
$1/N$ to the sum in Definition~\ref{def:sd} and each point of
$S_1 \setminus S_0$ contributes the same amount, and the factor $\tfrac12$
halves the total. Shifting an interval of length $2B'+1$ by $e$ leaves
$\abs{e}$ points uncovered at one end, so
$\SD(E, E+e) = \abs{e}/(2B'+1) \le B/(2\cdot 2^{\kappa}B) \le 2^{-\kappa}$. The
final claim follows from the triangle inequality applied through $E$.
\end{proof}

\begin{remark}
\label{rem:smudge-cost}
Lemma~\ref{lem:smudge} is cheap to state and expensive to use. It says that
hiding a $B$-bounded quantity statistically costs $\kappa$ \emph{additional} bits
of range in the mask, and in a lattice ciphertext that range is subtracted from
the noise budget that also has to accommodate the evaluation. This is the exact
mechanism by which the soundness slack of a proof about a ciphertext is paid for
out of the depth available for computing on it, and we report the measured
consequence in Section~\ref{sec:threats}.
\end{remark}

\subsection{Algebraic structures, and why the choice of ring is not cosmetic}
\label{sec:math:algebra}

Two ambient structures appear throughout: the ring $\Z_{2^k}$ of machine words
and the field $\mathbb{F}_p$ of integers modulo a prime. They are not
interchangeable, and the differences propagate all the way into our cost
measurements, so we set them out explicitly.

\begin{definition}
\label{def:ringfield}
A \ib{commutative ring} $(\Ring,+,\cdot)$ is an abelian group under $+$ with an
associative, commutative multiplication distributing over $+$ and possessing an
identity $1$. An element $u \in \Ring$ is a \ib{unit} when $uv = 1$ for some
$v$, and a \ib{zero divisor} when $u \ne 0$ and $uv = 0$ for some $v \ne 0$. A
\ib{field} is a commutative ring in which every nonzero element is a unit.
\end{definition}

\begin{proposition}
\label{prop:units}
In $\Z_m$, an element $a$ is a unit if and only if $\gcd(a,m)=1$. In particular
$\Z_{2^k}$ has exactly $2^{k-1}$ units, namely the odd residues, and every even
nonzero residue is a zero divisor; $\Z_p$ for $p$ prime is a field.
\end{proposition}

\begin{proof}
If $\gcd(a,m)=1$ then by B\'ezout there are $u,v$ with $au+mv=1$, so $au \equiv 1
\pmod m$. Conversely if $au \equiv 1 \pmod m$ then $au - 1 = mv$ for some $v$, so
any common divisor of $a$ and $m$ divides $1$. For $m = 2^k$ the condition
$\gcd(a,2^k)=1$ says exactly that $a$ is odd, and there are $2^{k-1}$ odd
residues; if $a = 2^s b$ with $1 \le s < k$ then $a \cdot 2^{k-s} \equiv 0$, so
$a$ is a zero divisor. For $m = p$ prime, every $a \in \{1,\dots,p-1\}$ is
coprime to $p$.
\end{proof}

Proposition~\ref{prop:units} has three consequences that recur in this survey,
and we state them together because each is a place where a textbook protocol has
to be modified before it can run over machine words.

\begin{itemize}
  \item \textbf{Shamir sharing does not transfer.} Reconstruction interpolates a
  polynomial by Lagrange's formula, whose coefficients are differences of
  evaluation points inverted. Over $\Z_{2^k}$ those differences are frequently
  even and therefore not invertible, so the scheme is not well defined for
  arbitrary evaluation points. This is why our ring-based layer uses
  \emph{replicated} instead of threshold sharing
  (Definition~\ref{def:sharings}), and why the field-based engines in our
  compiler use a Mersenne prime instead.
  \item \textbf{Statistical masking must respect the power-of-two structure.} A
  mask uniform on $\Z_{2^k}$ hides perfectly (Theorem~\ref{thm:mask}), but a mask
  uniform on a subinterval does not, and the failure is not visible to any
  correctness test, since the arithmetic is unaffected. An instance is recorded in
  Appendix~\ref{sec:appendix:caught}.
  \item \textbf{Truncation is natural.} Division by $2^d$ with rounding is a shift
  on a machine word, and fixed-point arithmetic wants precisely that operation.
  Over $\mathbb{F}_p$ the corresponding operation needs a comparison against
  $p/2$ to recover the sign, which is why the field engines pay for a wider
  statistical margin. Section~\ref{sec:math:fixed} makes the trade quantitative.
\end{itemize}

\begin{theorem}[Chinese remainder]
\label{thm:crt}
Let $m = m_1 m_2 \cdots m_r$ with the $m_i$ pairwise coprime. Then the map
$x \mapsto (x \bmod m_1, \dots, x \bmod m_r)$ is a ring isomorphism
$\Z_m \cong \Z_{m_1} \times \cdots \times \Z_{m_r}$.
\end{theorem}

\begin{proof}
The map is a ring homomorphism because reduction modulo each $m_i$ is. It is
injective: if $x \equiv y$ modulo every $m_i$ then each $m_i$ divides $x-y$, and
since the $m_i$ are pairwise coprime their product does, so $x = y$ in $\Z_m$.
Both sides have cardinality $m$, so an injective map between them is bijective.
\end{proof}

Theorem~\ref{thm:crt} is used in two distinct ways below, and conflating them is a
common source of confusion. In the \emph{implementation} of lattice arithmetic it
is a representation device: a large modulus $q$ is stored as a residue system over
machine-word primes, which is what makes a number-theoretic transform practical
(Section~\ref{sec:math:ntt}). In the \emph{plaintext} space of a homomorphic
scheme it is a parallelism device: the plaintext ring factors into independent
slots, so one ciphertext carries many independent values
(Section~\ref{sec:math:packing}). The first is invisible to the protocol; the
second changes its asymptotics.

\begin{definition}
\label{def:polyring}
For a commutative ring $\Ring$ we write $\Ring[X]$ for the ring of polynomials
with coefficients in $\Ring$. For $N \in \N$ the \ib{negacyclic} quotient is
$\Ring[X]/(X^N+1)$, and the \ib{cyclic} quotient is $\Ring[X]/(X^N-1)$. We write
$\Ring = \Z[X]/(X^N+1)$ and $\Rq = \Z_q[X]/(X^N+1)$ for the rings used by the
homomorphic schemes of Section~\ref{sec:prelim:fhe}, with $N$ a power of two.
\end{definition}

For $N$ a power of two, $X^N+1$ is the $2N$-th cyclotomic polynomial and is
irreducible over $\Q$, so $\Q[X]/(X^N+1)$ is a field and $\Ring$ is the ring of
integers of a cyclotomic number field. Two properties of the negacyclic quotient
are used constantly. First, reduction is free: $X^N \equiv -1$, so a product of
two degree-$(N-1)$ polynomials is reduced by subtracting the high half from the
low half not by long division. Second, multiplication is a \ib{negacyclic
convolution},
\begin{equation}
\label{eq:negconv}
  (a \cdot b)_k \;=\; \sum_{i+j=k} a_i b_j \;-\; \sum_{i+j=k+N} a_i b_j ,
\end{equation}
which is the identity that makes coefficient packing work
(Section~\ref{sec:math:packing}) and which our convolution benchmark exploits.

\subsection{Fast polynomial multiplication}
\label{sec:math:ntt}

Homomorphic evaluation is dominated by multiplication in $\Rq$, so the cost of
Equation~\eqref{eq:negconv} determines the cost of everything above it. Done
directly it is $\Theta(N^2)$ coefficient products. The number-theoretic transform
reduces it to $\Theta(N \log N)$, and since we report a $40.6\times$ speed-up
from implementing it, we give the construction.

\begin{definition}
\label{def:ntt}
Let $q$ be a prime with $q \equiv 1 \pmod{2N}$, let $\psi \in \Z_q$ be a
primitive $2N$-th root of unity, and put $\omega = \psi^2$, a primitive $N$-th
root. The \ib{negacyclic number-theoretic transform} of
$\mathbf{a} = (a_0,\dots,a_{N-1})$ is
\[
  \widehat{a}_j \;=\; \sum_{i=0}^{N-1} a_i\, \psi^{i}\, \omega^{ij}
  \qquad (0 \le j < N),
\]
that is, the $N$-point discrete transform applied to the sequence
$(a_i \psi^i)_i$.
\end{definition}

\begin{theorem}
\label{thm:ntt}
With the notation of Definition~\ref{def:ntt}, for all $a,b \in \Rq$ the
negacyclic product satisfies
$\widehat{(a\cdot b)}_j = \widehat{a}_j\,\widehat{b}_j$ for every $j$, and the
transform is invertible with
$a_i = N^{-1}\psi^{-i}\sum_j \widehat{a}_j \omega^{-ij}$.
Hence a negacyclic product costs two forward transforms, $N$ coefficient
products and one inverse transform.
\end{theorem}

\begin{proof}
Write $A(X) = \sum_i a_i X^i$ and let $c = a\cdot b$ in
$\Z_q[X]/(X^N+1)$. Because $\psi^N = -1$, the map $X \mapsto \psi$ sends
$X^N+1$ to $0$, so evaluation at $\psi\omega^{j}$ is a ring homomorphism from
$\Z_q[X]/(X^N+1)$ to $\Z_q$ for each $j$: indeed
$(\psi\omega^j)^N = \psi^N \omega^{jN} = -1$. Now
$\widehat{a}_j = \sum_i a_i \psi^i \omega^{ij} = A(\psi\omega^{j})$, so
$\widehat{c}_j = C(\psi\omega^j) = A(\psi\omega^j)B(\psi\omega^j) =
\widehat{a}_j\widehat{b}_j$, which is the first claim. The points
$\{\psi\omega^j\}_{j<N}$ are distinct and are exactly the $N$ roots of $X^N+1$ in
$\Z_q$, which exist because $q \equiv 1 \pmod {2N}$; a polynomial of degree less
than $N$ is determined by its values at $N$ distinct points, so the transform is
a bijection, and the stated inversion formula follows from the orthogonality
relation $\sum_{j} \omega^{j(i-i')} = N\,[i = i']$ in $\Z_q$.
\end{proof}

\begin{remark}
\label{rem:ntt-eng}
Theorem~\ref{thm:ntt} is a statement about $\Z_q$ for prime $q$. Deployed
parameters use a composite $q$, and the standard resolution is
Theorem~\ref{thm:crt}: represent $\Z_q$ as a residue number system over
machine-word primes each satisfying $q_i \equiv 1 \pmod{2N}$, and transform
independently in each. What remains after that is the cost of a single butterfly,
and the engineering lesson we report is that this cost is dominated by modular
reduction not by memory traffic, which is why a Shoup-preconditioned
multiply is worth more than a better loop order.
\end{remark}

\subsection{Fixed-point representation and its two error terms}
\label{sec:math:fixed}

Neither secret sharing over $\Z_{2^k}$ nor an integer homomorphic scheme offers
real arithmetic, and the networks we evaluate are specified over the reals. The
bridge is fixed point, and because two of our findings are about the cost of the
rescaling it forces, we set out its arithmetic precisely.

\begin{definition}
\label{def:fixed}
Fix a \ib{scale} $f \in \N$. The \ib{fixed-point encoding} of $x \in \Rr$ at
scale $f$ is $\lfloor x \cdot 2^{f} \rceil \in \Z$, and a stored integer $X$
decodes to $X \cdot 2^{-f}$. We say $X$ \ib{represents} $x$ with error
$\abs{X 2^{-f} - x} \le 2^{-f-1}$.
\end{definition}

Addition of encodings at a common scale is exact and needs no adjustment.
Multiplication does not preserve the scale: if $X$ and $Y$ encode $x$ and $y$ at
scale $f$ then $XY$ encodes $xy$ at scale $2f$, so a \ib{rescale} by $2^{-f}$ is
required before the value can be added to another scale-$f$ quantity. Two error
terms therefore accumulate, and they behave differently.

\begin{proposition}
\label{prop:fixed-error}
Let $X_i$ encode $x_i$ and $W_i$ encode $w_i$ at scale $f$, for
$i = 1,\dots,n$, with $\abs{x_i},\abs{w_i} \le M$. Let
$S = \rnd{2^{-f}\sum_i W_i X_i}$ be the rescaled inner product. Then
\[
  \abs{S\,2^{-f} - \textstyle\sum_i w_i x_i}
  \;\le\; \underbrace{n\,M\,2^{-f}}_{\text{input rounding}}
  \;+\; \underbrace{2^{-f-1}}_{\text{one rescale}}
  \;+\; O(n 2^{-2f}).
\]
\end{proposition}

\begin{proof}
Write $X_i = 2^f x_i + \delta_i$ and $W_i = 2^f w_i + \epsilon_i$ with
$\abs{\delta_i},\abs{\epsilon_i} \le \tfrac12$. Expanding the product pairs each
encoded value with the \emph{other} factor's rounding error,
\[
  W_i X_i \;=\; 2^{2f} w_i x_i \;+\; 2^{f}\bigl(w_i \delta_i + x_i \epsilon_i\bigr)
  \;+\; \delta_i \epsilon_i .
\]
Summing over $i$ and multiplying by $2^{-f}$,
\[
  2^{-f}\sum_i W_iX_i \;=\; 2^{f}\sum_i w_ix_i
  \;+\; \sum_i \bigl(w_i\delta_i + x_i\epsilon_i\bigr)
  \;+\; 2^{-f}\sum_i \delta_i\epsilon_i .
\]
The second term is at most $n M$ in absolute value, since
$\abs{w_i},\abs{x_i} \le M$ and each error is at most $\tfrac12$; the third is at
most $n2^{-f}/4$. The final rounding to an integer contributes at most
$\tfrac12$. Dividing throughout by $2^{f}$ gives the claim.
\end{proof}

Proposition~\ref{prop:fixed-error} is the reason the two paradigms diverge on
nonlinearity cost, and it is worth stating why explicitly. The dominant term is
input rounding, which is inherited from quantisation and is paid once per value
regardless of protocol. The rescale term is a single half-unit, negligible in
accuracy---and yet, in a secret-sharing protocol, the rescale is the expensive
operation, because a right shift of a \emph{secret} value is not a local
operation over $\Z_{2^k}$: the shares of $X$ do not shift independently, since
the sum of the shifted shares differs from the shift of the sum by the carries
lost at the boundary. Section~\ref{sec:math:bitdec} gives the protocol and its
round cost. In a homomorphic scheme, by contrast, the rescale is either free
(BFV/BGV at integer scale, where the plaintext modulus absorbs it) or a native
operation on the modulus chain (CKKS, where it is exactly modulus switching).
This asymmetry, and not the multiplication count, is what our measurements keep
finding.

\begin{remark}
\label{rem:x2}
An activation choice interacts with the same arithmetic. For a comparison-free
scheme the natural nonlinearity is $x \mapsto x^2$, which is exact in both
paradigms. But squaring a scale-$f$ encoding produces a scale-$2f$ value, so it
forces a rescale of its own, and in a secret-sharing protocol that rescale costs
exactly what a comparison would have cost. It also has a dynamic-range
consequence: if a layer's linear map has gain $g$, then the magnitude of the
activations obeys $M_{\ell+1} \approx g^2 M_\ell^2$, which for $g M > 1$ grows
doubly exponentially in depth and overflows any fixed word length after a few
layers. Scaling weights so that $gM \approx 1$---which is what standard
initialisation does---puts the recursion at its fixed point. We report both the
failure and the fix in Section~\ref{sec:matched}.
\end{remark}

\subsection{Communication complexity and the cost model}
\label{sec:math:cost}

Our comparisons are quoted in wall-clock time computed from measured counts, so
the model that performs that conversion is part of the foundations and not a detail of the harness.

\begin{definition}
\label{def:roundsbytes}
Fix a protocol execution. Its \ib{communication} is the total number of bytes a
single party sends, and its \ib{round complexity} is the length of the longest
chain of messages in which each depends on the previous one. Formally, orient the
message graph by causal dependence; the round complexity is the height of that
directed acyclic graph.
\end{definition}

\begin{definition}
\label{def:netmodel}
For a link with one-way latency $L$ and per-party bandwidth $B$, the
\ib{modelled latency} of an execution with $R$ rounds, $C$ bytes of
communication and $T$ seconds of local computation is
\[
  \mathrm{Time} \;=\; R \cdot L \;+\; C / B \;+\; T .
\]
\end{definition}

Definition~\ref{def:netmodel} is a lower bound made into an estimate. It is a
lower bound because a party cannot receive a causally dependent message sooner
than $L$ after sending its predecessor, and cannot transfer $C$ bytes faster than
$B$; it becomes an estimate by assuming perfect overlap within a round and no
overlap between rounds. We validate it against a socket to within
\SI{13.6}{\percent} worst case (Section~\ref{sec:method:net}), and we record in
Section~\ref{sec:threats} the one place where dropping $T$ made a local-area
comparison wrong by a factor of nineteen.

\begin{proposition}
\label{prop:asap}
Let a program be a directed acyclic graph of operations in which each operation
is either \ib{local} or \ib{interactive}, and suppose every set of interactive
operations that are mutually independent can be executed in a single round.
Schedule each operation at the earliest level permitted by its dependences. Then
the number of rounds equals the maximum number of interactive operations on any
directed path, and no schedule achieves fewer.
\end{proposition}

\begin{proof}
Write $d(v)$ for the maximum number of interactive operations on a path ending at
$v$, so that $d$ is computed by the recurrence $d(v) = \max_{u \to v} d(u) + [v
\text{ interactive}]$. Assign $v$ to level $d(v)$. Two interactive operations at
the same level cannot be connected by a path, since a path between them would
strictly increase $d$; hence each level's interactive operations are mutually
independent and by hypothesis occupy one round, giving $\max_v d(v)$ rounds.
Conversely, along a path carrying $m$ interactive operations, each must complete
before the next begins, so any schedule uses at least $m$ rounds.
\end{proof}

Proposition~\ref{prop:asap} is why we report \emph{dependence depth} not
operation counts as the predictor of latency, and why batching a layer's
comparisons is worth more than reducing their number: a layer of $n$ independent
comparisons costs one comparison's rounds. It also explains a structural feature
of the assignment problem in Section~\ref{sec:assign}: because a level costs the
maximum over its members not the sum, the objective is not separable
across operations, and an integer-linear formulation of it is not faithful.

\section{Foundations of the protocols}
\label{sec:found}

Section~\ref{sec:math} fixed the mathematics; this section builds the protocol
machinery on it. Section~\ref{sec:prelim} then presents the specific
constructions our artifact implements, in the form it implements them. The
division is that everything here holds for any implementation, whereas
Section~\ref{sec:prelim} records our choices.

\subsection{The real/ideal paradigm}
\label{sec:found:sim}

Definition~\ref{def:corruption} named the corruption models; we now give the
security definition itself, because our composability argument
(Theorem~\ref{thm:uc}) quantifies over it and a reader cannot check that argument
against an informal notion.

\begin{definition}
\label{def:functionality}
An \ib{ideal functionality} $\Fun$ for $n$ parties is a probabilistic
interactive machine that receives an input $x_i$ from each $P_i$, computes
$(y_1,\dots,y_n) \gets \Fun(x_1,\dots,x_n)$, and delivers $y_i$ to $P_i$. A
functionality may additionally be given a \ib{leakage} function $\mathcal{L}$,
whose value is delivered to the adversary; a functionality with no leakage
function is \ib{leakage-free}.
\end{definition}

Making leakage explicit in the functionality is not pedantry. A protocol that
opens an intermediate value is not insecure if the functionality says the value is
opened; it is insecure if the functionality does not. Our own broken truncation
(Section~\ref{sec:disclosure}) is exactly a case where the protocol opened
slightly more than the functionality we had written down, and stating the
functionality with its leakage clause is what made the discrepancy visible.

\begin{definition}
\label{def:realideal}
Let $\Prot$ be an $n$-party protocol and $\Fun$ a functionality. For an adversary
$\Adv$ corrupting a set $I \subseteq \{P_0,\dots,P_{n-1}\}$ with $\abs{I} \le t$,
let $\real_{\Prot,\Adv,I}(\mathbf{x},\lambda)$ denote the joint distribution of
$\Adv$'s view and the honest parties' outputs in an execution of $\Prot$ on
inputs $\mathbf{x}$. For a simulator $\Sm$, let
$\ideal_{\Fun,\Sm,I}(\mathbf{x},\lambda)$ denote the same joint distribution in
an ideal execution where the parties hand their inputs to $\Fun$ and $\Sm$ plays
the corrupted parties given only $\{x_i\}_{i \in I}$, $\{y_i\}_{i\in I}$ and any
leakage. We say $\Prot$ \ib{securely realises} $\Fun$ against $t$ static
corruptions when for every such $\Adv$ there exists a probabilistic
polynomial-time $\Sm$ with
\[
  \bigl\{\real_{\Prot,\Adv,I}(\mathbf{x},\lambda)\bigr\}
  \;\approx\;
  \bigl\{\ideal_{\Fun,\Sm,I}(\mathbf{x},\lambda)\bigr\}
\]
for all $\mathbf{x}$ and all admissible $I$, where $\approx$ is perfect,
statistical or computational indistinguishability as specified.
\end{definition}

Three features of Definition~\ref{def:realideal} matter to us. It is a
\emph{joint} distribution over the adversary's view and the honest outputs, so a
protocol cannot be excused for letting the adversary influence the result while
learning nothing. The simulator is given the corrupted parties' outputs, which is
what makes simulation possible at all: whatever the adversary sees must be
reconstructible from what it is entitled to know. And corruption here is
\ib{static}, fixed before the execution; \ib{adaptive} corruption, where the
adversary chooses whom to corrupt as it goes, is strictly harder and we do not
claim it.

\begin{remark}
\label{rem:god}
Definition~\ref{def:realideal} says nothing about whether honest parties obtain
output. Three increasingly strong guarantees are distinguished in the literature:
\ib{security with abort}, where honest parties may halt without output;
\ib{fairness}, where the adversary obtains output only if honest parties do; and
\ib{guaranteed output delivery}. With a dishonest majority the last two are
unachievable in general---Cleve's impossibility result for coin flipping applies
already to two parties---so a dishonest-majority protocol offering abort is not
making a weak choice but the only available one. Our artifact offers abort in an
honest-majority setting, where the stronger guarantees are achievable and we do
not achieve them; we say so rather than imply otherwise.
\end{remark}

\subsection{Secret sharing}
\label{sec:found:ss}

\begin{definition}
\label{def:ss}
Let $\Ring$ be a finite commutative ring and let $\mathcal{A}$ be a monotone
collection of subsets of $\{0,\dots,n-1\}$, the \ib{authorised sets}. A
\ib{secret-sharing scheme} is a pair of algorithms $(\mathsf{Share},
\mathsf{Rec})$ where $\mathsf{Share}(x)$ outputs $(s_0,\dots,s_{n-1})$ and
$\mathsf{Rec}$ maps the shares held by any authorised set back to $x$. It is
\ib{correct} when $\mathsf{Rec}$ succeeds for every authorised set, and
\ib{perfectly private} when for every unauthorised set $U \notin \mathcal{A}$ and
every pair of secrets $x,x'$, the distributions of $\{s_i\}_{i \in U}$ under
$\mathsf{Share}(x)$ and $\mathsf{Share}(x')$ are identical. It is \ib{linear}
when $\mathsf{Share}(x)+\mathsf{Share}(y)$ is a valid sharing of $x+y$ under
componentwise addition.
\end{definition}

Linearity is the property that makes secret sharing useful for computation:
addition and multiplication by a public constant are local, so only multiplication
of two secrets needs interaction, and the entire cost structure of the paradigm
follows from that split.

\begin{proposition}
\label{prop:additive-priv}
Additive sharing over $\Ring$, in which $s_0,\dots,s_{n-2}$ are uniform and
$s_{n-1} = x - \sum_{i<n-1} s_i$, is correct for the single authorised set
$\{0,\dots,n-1\}$ and perfectly private against every proper subset.
\end{proposition}

\begin{proof}
Correctness is immediate. For privacy, let $U$ be a proper subset, so that
$U$ omits at least one index. There are two cases. If $n-1 \notin U$, then
$\{s_i\}_{i \in U}$ is a subset of the independent uniform values
$s_0,\dots,s_{n-2}$, whose joint distribution does not involve $x$ at all. If
$n-1 \in U$, pick $j \le n-2$ with $j \notin U$, which exists because $U$ is
proper. Condition on $\{s_i\}_{i \in U \setminus \{n-1\}}$, all of which are
uniform and independent of $x$. Then
$s_{n-1} = x - \sum_{i \le n-2} s_i$ contains the term $-s_j$ with $s_j$ uniform,
unobserved and independent of everything conditioned on, so $s_{n-1}$ is uniform
on $\Ring$ and independent of $x$. In both cases the observed tuple is
distributed identically for every secret. We remark that this argument uses only
that $\Ring$ is an abelian group under addition, so it holds over $\Z_{2^k}$
despite Proposition~\ref{prop:units}---which is exactly why additive and
replicated sharing transfer to machine words while Shamir sharing does not.
\end{proof}

\begin{proposition}
\label{prop:rep-priv}
Replicated sharing among three parties (Definition~\ref{def:sharings}), in which
$P_i$ holds $(x_i,x_{i+1})$ for an additive sharing $x_0+x_1+x_2 = x$, is correct
for any two parties and perfectly private against any one.
\end{proposition}

\begin{proof}
Any two parties jointly hold all three additive shares: $P_i$ and $P_{i+1}$ hold
$\{x_i,x_{i+1}\}$ and $\{x_{i+1},x_{i+2}\}$, whose union is all of them, so
reconstruction succeeds. A single party $P_i$ holds $(x_i,x_{i+1})$ and is missing
$x_{i+2}$, which is uniform and independent of the pair it holds; by the argument
of Proposition~\ref{prop:additive-priv} the observed pair is identically
distributed for every secret.
\end{proof}

\begin{definition}
\label{def:shamir}
Let $\mathbb{F}$ be a finite field with $\abs{\mathbb{F}} > n$ and fix distinct
nonzero evaluation points $\alpha_0,\dots,\alpha_{n-1}$. \ib{Shamir sharing} of
$x$ with threshold $t$ samples a uniform polynomial $p$ of degree at most $t$
subject to $p(0)=x$ and sets $s_i = p(\alpha_i)$. The authorised sets are those
of size at least $t+1$.
\end{definition}

\begin{proposition}
\label{prop:shamir}
Shamir sharing is correct and perfectly private, and it is linear.
\end{proposition}

\begin{proof}
Correctness is Lagrange interpolation: given $t+1$ points, the unique polynomial
of degree at most $t$ through them is
$p(X) = \sum_{j} s_{i_j} \prod_{l \ne j} \frac{X - \alpha_{i_l}}{\alpha_{i_j} -
\alpha_{i_l}}$, and $x = p(0)$. This is where a field is needed: the denominators
are differences of distinct evaluation points and must be invertible, which
Proposition~\ref{prop:units} shows can fail over $\Z_{2^k}$. For privacy, fix an
unauthorised set of size $t$ and any candidate secret $x'$; the $t$ observed
values together with $p(0)=x'$ determine $t+1$ conditions on a polynomial with
$t+1$ free coefficients, so exactly one polynomial is consistent, and the map
from the sampling randomness to the observed shares is a bijection for every
$x'$. Linearity holds because the sum of two polynomials of degree at most $t$
has degree at most $t$ and evaluates pointwise.
\end{proof}

\begin{remark}
\label{rem:degree}
Multiplying two Shamir sharings pointwise yields a sharing of the product on a
polynomial of degree up to $2t$, which is not a valid degree-$t$ sharing. The
standard repair, \ib{degree reduction}, has each party reshare its product and
the parties take a fixed linear combination; it costs one round and is why a
Shamir multiplication is interactive even with an honest majority. Replicated
sharing avoids the degree problem entirely at the cost of storing two shares per
party, which is the trade Definition~\ref{def:sharings} makes.
\end{remark}

\subsection{Multiplication: from correlated randomness, and from triples}
\label{sec:found:mul}

\begin{definition}
\label{def:beaver}
A \ib{Beaver triple} over $\Ring$ is a sharing $(\shr{a},\shr{b},\shr{c})$ with
$a,b$ uniform and $c = ab$.
\end{definition}

\begin{theorem}
\label{thm:beaver}
Given a Beaver triple and sharings $\shr{x},\shr{y}$, the protocol that opens
$d = x - a$ and $e = y - b$ and outputs
$\shr{z} = \shr{c} + d\shr{b} + e\shr{a} + de$
is correct, uses one round and two openings, and leaks nothing: the opened values
are uniform and independent of $x,y$.
\end{theorem}

\begin{proof}
For correctness, expand using $c = ab$:
\[
  c + db + ea + de = ab + (x-a)b + (y-b)a + (x-a)(y-b)
  = ab + xb - ab + ya - ab + xy - xb - ya + ab = xy .
\]
Each term on the left is either a sharing scaled by a public value or a public
constant, so the combination is local and the result is a valid sharing by
linearity. For privacy, $a$ is uniform and independent of $x$, so $d = x-a$ is
uniform on $\Ring$ and independent of $x$; likewise $e$. A simulator therefore
samples $d,e$ uniformly, which is exactly the real distribution, and the
simulation is perfect. Both openings are independent and occupy one round.
\end{proof}

Theorem~\ref{thm:beaver} explains the shape of the dishonest-majority cost we
measure: the online phase is two openings and is cheap, and the entire expense is
producing $(\shr{a},\shr{b},\shr{c})$, which needs either oblivious transfer or a
homomorphic scheme. Section~\ref{sec:prelim:mpc} gives the honest-majority
alternative, Theorem~\ref{thm:mul}, in which correlated randomness from a pairwise
pseudorandom function replaces the triple and the multiplication needs no
preprocessing at all. That contrast---preprocessing volume against threat
model---is the mechanism behind the crossover of
Section~\ref{sec:crossover-fhe}.

\begin{proposition}
\label{prop:bits-from-mul}
Let each party $P_i$ contribute a uniform bit $b_i \in \{0,1\}$ as a private
input. Then $b = \bigoplus_i b_i$ is a uniform bit unknown to any proper subset
of the parties, and it can be computed from the $\shr{b_i}$ by
$\ceil{\log n}$ rounds of multiplication using the identity
$u \oplus v = u + v - 2uv$ over $\Ring$ with $u,v \in \{0,1\}$.
\end{proposition}

\begin{proof}
The identity holds by inspection on the four cases. Uniformity: conditioned on
any proper subset of the $b_i$, at least one $b_j$ remains uniform and
independent, and exclusive-or with a uniform independent bit yields a uniform
bit. Balanced binary reduction of $n$ terms has depth $\ceil{\log n}$, and all
exclusive-ors at one level are independent, so by Proposition~\ref{prop:asap}
each level costs one round.
\end{proof}

Proposition~\ref{prop:bits-from-mul} is how our engines obtain shared random bits
without a trusted dealer, and it prices them: one multiplication per bit per
level, hence one Beaver triple per bit at two parties. We record in
Section~\ref{sec:disclosure} that treating these bits as free---the usual
``fake offline'' simplification---understated the cost of a comparison by a large
factor, because a comparison consumes tens of them.

\subsection{Comparison, bit decomposition and truncation}
\label{sec:math:bitdec}

Every nonlinearity we evaluate reduces to the sign of a shared value:
$\ReLU(x) = x \cdot [x > 0]$, a maximum is a mux on a comparison, and an argmax
is a fold of maxima. Over $\Z_{2^k}$ the sign is the top bit, so the primitive is
bit decomposition, and its cost dominates our secret-sharing measurements.

\begin{definition}
\label{def:bitdec}
Given $\shr{x}$ with $x$ known to lie in $[-2^{\ell-1},2^{\ell-1})$, a
\ib{bit-decomposition} protocol outputs sharings $\shr{x_0},\dots,\shr{x_{\ell-1}}$
of the two's-complement bits of $x$.
\end{definition}

The construction we use is the masked-opening one. Draw $\ell + \kappa$ shared
random bits $r_j$ (Proposition~\ref{prop:bits-from-mul}), form
$\shr{r} = \sum_j 2^j \shr{r_j}$ locally, open
$c = x + 2^{\ell-1} + r$, and recover the bits of $u = x + 2^{\ell-1}$ as the
binary difference $c - r$ computed on the shared bits of $r$ against the public
bits of $c$. Two facts make this work, and we prove both because the constants are
ones we quote.

\begin{lemma}
\label{lem:mask-bits}
With $r = \sum_{j<\ell+\kappa} 2^j r_j$ for independent uniform bits $r_j$, and
$u \in [0,2^{\ell})$, the opened value $c = u + r$ has statistical distance at most
$2^{-\kappa}$ from a distribution independent of $u$.
\end{lemma}

\begin{proof}
$r$ is uniform on $[0,2^{\ell+\kappa})$. By the argument of
Lemma~\ref{lem:smudge} applied to integer intervals, shifting a uniform interval
of cardinality $2^{\ell+\kappa}$ by $u \le 2^{\ell}-1$ leaves at most $u$ points
uncovered at each end, so
$\SD(r, r+u) \le u/2^{\ell+\kappa} < 2^{\ell}/2^{\ell+\kappa} = 2^{-\kappa}$.
Taking $T$ as in Definition~\ref{def:sd} gives the claim for any distinguisher.
\end{proof}

\begin{remark}
\label{rem:kappa-cost}
Lemma~\ref{lem:mask-bits} is why a comparison at $\ell$ bits consumes
$\ell + \kappa$ random bits rather than $\ell$, and therefore why the arithmetic
domain must be at least $\ell + \kappa + 1$ bits wide: the masked sum must not
wrap. With $\ell = 33$ and $\kappa = 26$ this is $60$ bits, which is precisely the
constraint that forces the field of our somewhat-homomorphic preprocessing engine
to be $61$ bits and hence forces its ciphertext modulus to be large. A reader
tracking why those parameters are what they are should start here.
\end{remark}

\begin{theorem}[Parallel-prefix borrow]
\label{thm:kogge}
Let $c$ be public and let the bits $r_j$ of $r$ be shared. The bits of
$u = c - r$ can be computed in $\ceil{\log \ell}$ rounds of shared
multiplication, independently of $\ell$'s value beyond that logarithm.
\end{theorem}

\begin{proof}
Write the borrow-out of position $j$ as $\beta_j$, with $\beta_{-1}=0$ and
\[
  \beta_j \;=\; g_j \;\vee\; (p_j \wedge \beta_{j-1}), \qquad
  g_j = \lnot c_j \wedge r_j, \qquad p_j = \lnot(c_j \oplus r_j),
\]
and the difference bit $u_j = c_j \oplus r_j \oplus \beta_{j-1}$. Because $c$ is
\emph{public}, $g_j$ and $p_j$ are local: each collapses to $r_j$ or $1-r_j$
according to the public bit $c_j$, so no interaction is needed to form them.
Define the associative operator
$(g,p)\circ(g',p') = (g \vee (p \wedge g'),\; p \wedge p')$, whose associativity
is checked by expanding both bracketings. Then $\beta_j$ is the first component
of $(g_j,p_j)\circ\cdots\circ(g_0,p_0)$, and a balanced prefix network
(Kogge--Stone) computes all $\ell$ prefixes in $\ceil{\log \ell}$ levels. It
remains to observe that $g_j p_j = 0$ always---if $c_j = 0$ then $g_j = r_j$ and
$p_j = 1-r_j$, and if $c_j = 1$ then $g_j = 0$---so the disjunction in $\circ$ may
be replaced by addition, making each application two shared multiplications. All
applications at one level are independent, so by Proposition~\ref{prop:asap} each
level is one round.
\end{proof}

\begin{corolary}
\label{cor:cmp-cost}
A comparison of two $\ell$-bit shared values costs one opening, $\ell+\kappa$
shared random bits, and $\ceil{\log(\ell+1)}$ rounds of multiplication; a batch
of any number of independent comparisons costs the same number of rounds as one.
\end{corolary}

\begin{proof}
Compare $x$ against $y$ by decomposing $x-y$ and taking the sign bit, which is
the top bit of the two's-complement representation; the counts are those of
Lemma~\ref{lem:mask-bits} and Theorem~\ref{thm:kogge} with $\ell+1$ bits to
accommodate the difference. The batching claim is Proposition~\ref{prop:asap}
applied to a set of mutually independent decompositions.
\end{proof}

Corollary~\ref{cor:cmp-cost} is the single most consequential statement in this
section for our measurements. It says that the nonlinear surface of a network
costs \emph{rounds proportional to depth and not to width}, which is why our
latency figures scale with the number of layers and are almost independent of the
number of neurons per layer, and why the dominant term in a comparison's cost is
the random bits instead of the prefix network.

\begin{proposition}[Truncation]
\label{prop:trunc}
Let $\shr{x}$ be given with the bits of $x$ available as $\shr{x_j}$, and let
$d < \ell$. Then $\shr{\floorr{x/2^{d}}}$ is obtained locally as
\[
  \shr{\floorr{x/2^{d}}} \;=\; \sum_{j=d}^{\ell-1} 2^{\,j-d}\,\shr{x_j}
  \;-\; 2^{\,\ell-d}\,\shr{x_{\ell-1}} ,
\]
that is, by recomposing the high bits and then correcting the sign; no
interaction is required beyond the decomposition.
\end{proposition}

\begin{proof}
In two's complement over $\ell$ bits,
$x = -2^{\ell-1}x_{\ell-1} + \sum_{j<\ell-1}2^j x_j$. An arithmetic shift right by
$d$ is floor division by $2^d$, and
$\floorr{x/2^d} = -2^{\ell-1-d}x_{\ell-1} + \sum_{d \le j < \ell-1} 2^{j-d}x_j$,
because discarding the low $d$ bits of a two's-complement representation is
exactly flooring. The displayed expression recomposes
$\sum_{d \le j < \ell} 2^{j-d} x_j$, which includes $+2^{\ell-1-d}x_{\ell-1}$,
and then subtracts $2^{\ell-d} x_{\ell-1} = 2\cdot 2^{\ell-1-d}x_{\ell-1}$, which
converts that term to $-2^{\ell-1-d}x_{\ell-1}$. All operations are additions and public scalings, so
they are local.
\end{proof}

\begin{remark}
\label{rem:trunc-sign}
The sign correction in Proposition~\ref{prop:trunc} is easy to omit, and omitting
it yields a \emph{logical} and not an arithmetic shift: the result is correct for
non-negative inputs and wrong for every negative one. Since roughly half of a
network's pre-activations are negative, the error is not subtle in effect---but it
is invisible to any test whose fixtures are non-negative. This is a case where the
natural test set and the failure set are disjoint, which is why the correctness
condition is worth stating as a proposition not left to the
implementation.
\end{remark}

\subsection{Oblivious transfer and its extensions}
\label{sec:found:ot}

Dishonest-majority preprocessing rests on oblivious transfer, so we define it and
give the two constructions whose costs we measure.

\begin{definition}
\label{def:ot}
The functionality $\Fun_{\mathrm{OT}}$ takes messages $(m_0,m_1)$ from a sender
and a bit $b$ from a receiver, and delivers $m_b$ to the receiver and nothing to
the sender. In \ib{random OT} the sender has no input and receives uniform
$(m_0,m_1)$ as output; in \ib{correlated OT} with correlation $\Delta$ the sender
inputs $\Delta$, receives uniform $m_0$, and the receiver obtains
$m_b = m_0 + b\Delta$.
\end{definition}

\begin{proposition}[Gilboa]
\label{prop:gilboa}
Let $P_0$ hold $x \in \Z_{2^k}$ and $P_1$ hold $y \in \Z_{2^k}$. Using $k$
correlated OTs in which $P_0$ is the receiver with choice bits the binary digits
of $x$ and $P_1$ offers correlation $2^{i}y$ on the $i$-th transfer, the parties
obtain an additive sharing of $xy$.
\end{proposition}

\begin{proof}
On transfer $i$, $P_1$ holds a uniform $r_i$ and $P_0$ obtains
$r_i + x_i 2^{i} y$ where $x_i$ is the $i$-th bit of $x$. Summing,
$P_0$ holds $\sum_i r_i + y\sum_i x_i 2^i = \sum_i r_i + xy$ and $P_1$ holds
$-\sum_i r_i$; the two sum to $xy$. Privacy is inherited from the OT: $P_0$ sees
only its selected outputs, which are uniform given the unselected ones, and $P_1$
learns nothing about the choice bits.
\end{proof}

Proposition~\ref{prop:gilboa} prices a Beaver triple over $\Z_{2^k}$ at $2k$
correlated OTs, since a triple needs the two cross terms $a_0b_1$ and $a_1b_0$.
With $k=64$ and a per-OT cost of $\kappa$ bits plus a correction, this is the
$\sim\!3$ kilobytes per triple that our measurements confirm to within one per
cent, and it is the quantity against which somewhat-homomorphic preprocessing
competes.

\begin{remark}
\label{rem:iknp}
Base oblivious transfer costs public-key operations, so a protocol needing
millions of them uses an \ib{extension}: a small number of base OTs are turned
into arbitrarily many by symmetric-key work, at the cost of a correlation the
receiver's choices induce across the outputs. The construction transposes a
matrix of pseudorandom columns so that the sender learns, for each row, one of
two values determined by a global secret $\Delta$, and hashing breaks the
correlation. A malicious receiver can use inconsistent choice bits across
columns, which leaks bits of $\Delta$; a consistency check based on a random
linear combination over $\mathbb{F}_{2^{128}}$ closes this at $O(1)$ additional
communication. Our layer implements the extension and the check, and we report
that the extension is $27\times$ cheaper per transfer than the base protocol at
batch size $8192$.
\end{remark}

\subsection{Garbling schemes}
\label{sec:found:gc}

Definitions~\ref{def:garbling}--\ref{def:freexor} and
Theorem~\ref{thm:halfgates} give the constructions we implement; here we give the
security notions they must satisfy, because the cheapest optimisations are
exactly the ones that strengthen the assumption.

\begin{definition}
\label{def:gc-security}
A garbling scheme $(\Gb,\En,\Ev,\De)$ is \ib{correct} when
$\De(d,\Ev(F,\En(e,x))) = f(x)$ for all $x$. It is \ib{private} when there is a
simulator that, given $f$ and $f(x)$, produces $(F,\En(e,x),d)$ indistinguishable
from a real garbling; \ib{oblivious} when a simulator given only $f$ produces
$(F,\En(e,x))$ indistinguishably, that is, the evaluator learns not even the
output; and \ib{authentic} when no adversary given $(F,\En(e,x))$ can produce an
output label that decodes to anything other than $f(x)$.
\end{definition}

Privacy is what a semi-honest evaluator requires. Authenticity is what makes a
malicious evaluator unable to forge a result, and it is why the garbled-circuit
route can be given active security cheaply on the evaluator's side. Obliviousness
matters in our setting because the garbled output is kept as a sharing not
decoded, which is exactly how a garbled comparison composes with an arithmetic
protocol.

\begin{remark}
\label{rem:freexor-assumption}
Free exclusive-or replaces a garbled table by a wire-label offset, so that an
XOR gate costs nothing. The price is the assumption: security of the scheme now
requires the hash to be \ib{circular correlation robust}, because labels related
by the global offset $R$ are hashed together, whereas the plain construction needs
only a pseudorandom function. Half-gates then reduce an AND gate from three
ciphertexts to two while remaining compatible with free exclusive-or. The lesson
we take from this for our cost comparison is that the garbled route's headline
figure of two ciphertexts per AND is not assumption-free, and a fair comparison
against an arithmetic protocol whose security is information-theoretic should say
so.
\end{remark}

\subsection{Lattices, LWE and the hardness assumptions}
\label{sec:found:lattice}

Definition~\ref{def:lattice} introduces lattices and Definition~\ref{def:rlwe}
the ring problem our schemes rest on. Here we state the assumptions in the form
that parameter selection uses, since our security claims about the homomorphic
layer are exactly claims about these.

\begin{definition}
\label{def:lwe}
Let $n,q \in \N$ and let $\chi$ be a distribution over $\Z$ concentrated on small
values. The \ib{search LWE} problem is: given samples
$(\mathbf{a}_i, \langle \mathbf{a}_i,\mathbf{s}\rangle + e_i \bmod q)$ with
$\mathbf{a}_i \gets \Z_q^n$, $e_i \gets \chi$ and a fixed secret
$\mathbf{s} \in \Z_q^n$, recover $\mathbf{s}$. The \ib{decision LWE} problem is
to distinguish such samples from uniform samples over $\Z_q^n \times \Z_q$.
\end{definition}

\begin{definition}
\label{def:rlwe-restated}
The \ib{ring} variant replaces $\Z_q^n$ by $\Rq = \Z_q[X]/(X^N+1)$: samples are
$(a_i, a_i s + e_i)$ with $a_i \gets \Rq$, $s \in \Rq$ fixed and $e_i$ drawn
coefficientwise from $\chi$. A single ring sample carries $N$ scalar samples, so
the public key is smaller by a factor of $N$ and multiplication becomes a ring
product instead of a matrix--vector product.
\end{definition}

The assumptions we rely on are that decision LWE and decision ring-LWE are hard
for suitable parameters. Two facts about them shape our practice and not our proofs. First, search-to-decision and worst-case-to-average-case reductions are
known, which is why these problems are used at all; we cite not reprove
them. Second, and this is what a parameter table actually depends on, the
\emph{concrete} hardness is estimated by running the best known attacks---lattice
reduction with a sieving or enumeration subroutine, plus algebraic and hybrid
attacks---and reading off the estimated cost. Concrete security is therefore a
moving target, and the honest way to state a parameter set is to name the estimator
and its version. We do that in Section~\ref{sec:prelim:fhe} and we do not claim
security beyond what the standard tables assert.

\begin{remark}
\label{rem:ntru-lesson}
The distinction between ``a problem is hard'' and ``these parameters are hard'' is
not academic for us. We implemented a scheme whose security rested on a variant
assumption, and the variant admitted a subfield-lattice attack for the parameter
shape we had chosen; the construction was removed not reparameterised. We
record such a case in Appendix~\ref{sec:appendix:caught}. A survey that lists only
the primitives that survived would give a misleading picture of how the choice
between them is made.
\end{remark}

\subsection{Homomorphic encryption: noise, depth and the operations that manage them}
\label{sec:found:fhe}

Definition~\ref{def:bfvbgv} gives the two integer schemes we implement. This
subsection derives the noise arithmetic, because every structural claim we make
about homomorphic evaluation---why depth is bounded, why a bootstrap is needed,
why a scheme switch can replace one---is a statement about the growth of a single
quantity.

Let $\mathsf{ct} = (c_0,c_1) \in \Rq^2$ be an encryption of $m$ under secret $s$,
with the invariant, for the BFV encoding,
\begin{equation}
\label{eq:bfv-inv}
  c_0 + c_1 s \;=\; \Delta m + e \pmod q, \qquad \Delta = \floorr{q/t},
\end{equation}
where $t$ is the plaintext modulus and $e$ is the noise. Decryption computes
$c_0 + c_1 s$, divides by $\Delta$ and rounds.

\begin{proposition}
\label{prop:bfv-dec}
Decryption of Equation~\eqref{eq:bfv-inv} is correct provided
$\norm{e}_\infty < \Delta/2$.
\end{proposition}

\begin{proof}
We have $c_0 + c_1 s = \Delta m + e$ over $\Rq$, and lifting to centred
representatives, $\rnd{(\Delta m + e)/\Delta} = m + \rnd{e/\Delta} = m$ whenever
every coefficient of $e$ has magnitude below $\Delta/2$, since then
$\abs{e_i/\Delta} < 1/2$ and rounds to zero. If any coefficient exceeds the bound
the corresponding plaintext coefficient is off by at least one, and because
plaintexts are taken modulo $t$ the error is not detectable from the ciphertext.
\end{proof}

Proposition~\ref{prop:bfv-dec} is the entire budget. Everything homomorphic
evaluation does either consumes it or restores it.

\begin{proposition}
\label{prop:noise-growth}
Let $\mathsf{ct}_1,\mathsf{ct}_2$ encrypt $m_1,m_2$ with noise bounded by
$B_1,B_2$. Then
\begin{enumerate}
  \item addition yields noise at most $B_1+B_2$;
  \item multiplication by a plaintext of infinity-norm $\gamma$ and Hamming-weight
  at most $h$ yields noise at most $\gamma h B_1$;
  \item ciphertext multiplication yields a degree-two ciphertext whose noise is
  bounded by approximately $t N (B_1 + B_2) + \text{(terms independent of } B_i)$,
  that is, it grows \emph{multiplicatively} in the plaintext modulus and the ring
  dimension.
\end{enumerate}
\end{proposition}

\begin{proof}
For (1), add the invariants: $(c_0+c_0') + (c_1+c_1')s = \Delta(m_1+m_2) + (e_1+e_2)$,
and the triangle inequality bounds the sum. For (2), multiply the invariant by the
plaintext $p$; each output coefficient is a sum of at most $h$ products of a
coefficient of $p$ with one of $e_1$, whence the bound $\gamma h B_1$ by
Equation~\eqref{eq:negconv} and the triangle inequality. For (3), multiply the two
invariants:
\[
  (c_0 + c_1 s)(c'_0 + c'_1 s) = \Delta^2 m_1 m_2 + \Delta(m_1 e_2 + m_2 e_1) +
  e_1 e_2 ,
\]
whose left side expands to a quadratic in $s$, that is, a degree-two ciphertext
$(d_0,d_1,d_2)$ with $d_0 + d_1 s + d_2 s^2$ equal to the right side. Dividing by
$\Delta$ to restore the invariant scale introduces the factor $t$, and the
convolution of two degree-$(N-1)$ polynomials contributes up to $N$ terms per
coefficient, giving the stated shape. The dominant term is
$\Delta(m_1e_2 + m_2e_1)/\Delta = m_1e_2+m_2e_1$, bounded by $tN(B_1+B_2)/2$ since
plaintext coefficients are bounded by $t/2$.
\end{proof}

Three consequences follow, and they are the structural facts our measurements
keep meeting.

\begin{corolary}
\label{cor:depth}
With fresh noise $B_0$ and the multiplicative growth of
Proposition~\ref{prop:noise-growth}(3), the number of sequential ciphertext
multiplications before decryption fails---the \ib{multiplicative depth}---is
$O(\log_{tN}(\Delta/B_0))$, and is therefore \emph{logarithmic} in the ciphertext
modulus and \emph{linear} in $\log q$ for fixed $t,N$.
\end{corolary}

\begin{proof}
Each multiplication multiplies the noise bound by approximately $tN$, so after $L$
of them the bound is roughly $B_0 (tN)^{L}$. Decryption requires this to stay
below $\Delta/2 \approx q/2t$, giving $L \le \log_{tN}(q/(2tB_0))$.
\end{proof}

\begin{remark}
\label{rem:relin}
A ciphertext multiplication raises the degree in $s$ from one to two, so a second
multiplication would raise it to four and decryption would need $s^4$.
\ib{Relinearisation} restores degree one using an evaluation key that encrypts
$s^2$ under $s$; it costs a key-switching operation and adds noise. A scheme that
declines to relinearise is confined to \emph{one} ciphertext multiplication, and
that is exactly the regime our depth-one homomorphic layer occupies: it has no
evaluation keys at all, which removes their generation, their storage and their
noise, at the price of being unable to multiply twice. Whether that price is
payable is the question our scheme-switching result answers.
\end{remark}

\begin{remark}
\label{rem:modswitch}
\ib{Modulus switching} rescales a ciphertext from modulus $q$ to $q' < q$,
dividing both the ciphertext and its noise by approximately $q/q'$. It does not
improve the \emph{ratio} of noise to modulus by itself, but in BGV it is used to
keep that ratio constant after each multiplication, which converts multiplicative
noise growth into a linear consumption of a precomputed \ib{modulus chain}
$q_L > q_{L-1} > \cdots > q_0$. The number of links is the depth budget, and this
is why a levelled scheme's parameters are chosen from the depth of the circuit
before evaluation begins.
\end{remark}

\begin{definition}
\label{def:bootstrap}
\ib{Bootstrapping} is homomorphic evaluation of the decryption circuit, taking a
ciphertext whose noise is near the bound to a fresh ciphertext of the same
plaintext with noise near $B_0$. A scheme with bootstrapping is \ib{fully}
homomorphic; one without is \ib{levelled}.
\end{definition}

Bootstrapping is what makes unbounded depth possible and is, by a wide margin, the
most expensive operation in the paradigm: it requires evaluating a nontrivial
arithmetic circuit under the encryption it is repairing, at parameters chosen to
accommodate that circuit. The measurements we cite from the literature put a
single CKKS bootstrap at seconds to tens of seconds on a CPU, dominating the
inference that contains it. That cost, and the fact that a levelled pipeline which
declines to pay it simply stops at its depth bound, is the discontinuity our
comparison in Section~\ref{sec:matched} is organised around.

\begin{remark}
\label{rem:switch-as-bootstrap}
Definition~\ref{def:bootstrap} refreshes noise by computing decryption
homomorphically. A scheme switch achieves the same effect by a different route:
the ciphertext is decrypted \emph{under a sharing}, so the plaintext is never
exposed, and re-encrypted fresh. The cost is a round of interaction with parties
who hold the key in shares, instead of a homomorphic evaluation of decryption.
Whether this is a better trade is precisely a measurement question, and it is the
one our hybrid result answers: the switch is milliseconds where the bootstrap is
seconds, and the price is that the parties must be online.
\end{remark}

\subsection{Packing: two encodings and the operations they admit}
\label{sec:math:packing}

A ciphertext holds a ring element with $N$ coefficients, and using it to carry one
scalar wastes a factor of $N$. There are two ways to fill it, they support
different operations, and choosing between them changes the asymptotics of a
layer, so we set them out together.

\begin{definition}
\label{def:coefpack}
\ib{Coefficient packing} encodes a vector $\mathbf{a} \in \Z_t^N$ as the
polynomial $\sum_i a_i X^i$. \ib{Slot packing} fixes a plaintext modulus $t$ with
$t \equiv 1 \pmod{2N}$, so that $X^N+1$ splits into distinct linear factors over
$\Z_t$, and encodes $\mathbf{a}$ as the unique polynomial with
$a(\zeta_j) = a_j$ at the $N$ roots $\zeta_j$; by Theorem~\ref{thm:crt} this is a
ring isomorphism $\Z_t[X]/(X^N+1) \cong \Z_t^N$.
\end{definition}

\begin{proposition}
\label{prop:packing-ops}
Under coefficient packing, the product of two encodings is the negacyclic
convolution of the vectors, Equation~\eqref{eq:negconv}. Under slot packing, the
product is the componentwise product. Consequently a single ciphertext
multiplication computes, respectively, all $N$ shifted inner products of two
vectors, or $N$ independent scalar products.
\end{proposition}

\begin{proof}
The first claim is Definition~\ref{def:polyring} and
Equation~\eqref{eq:negconv}. The second is that evaluation at each $\zeta_j$ is a
ring homomorphism, so $(ab)(\zeta_j) = a(\zeta_j)b(\zeta_j)$, which is exactly
componentwise multiplication of the slot vectors; the isomorphism of
Theorem~\ref{thm:crt} makes the correspondence bijective.
\end{proof}

Proposition~\ref{prop:packing-ops} is the reason the two encodings suit different
layers. A convolution is a shifted inner product, so coefficient packing computes
a whole convolutional layer in one ciphertext multiplication with no rotations at
all: the $k^2$ taps of a kernel land in the coefficients of a single product, and
the desired outputs are read off at a stride. A matrix--vector product, by
contrast, needs sums \emph{across} positions, which coefficient packing does not
provide and slot packing provides only through rotations.

\begin{definition}
\label{def:rotation}
A \ib{rotation} by $k$ maps the slot vector $(a_0,\dots,a_{N-1})$ to its cyclic
shift. Homomorphically it is realised by the Galois automorphism
$X \mapsto X^{g}$ for a suitable $g$, followed by a key switch under a
\ib{Galois key}, and costs one key-switching operation.
\end{definition}

\begin{proposition}
\label{prop:hs-matvec}
A matrix--vector product with an $N \times N$ plaintext matrix on a slot-packed
ciphertext can be computed with $N$ plaintext multiplications and $N-1$ rotations
by summing the \ib{diagonals} of the matrix, and with
$O(\sqrt{N})$ rotations by the baby-step giant-step decomposition.
\end{proposition}

\begin{proof}
Write $\mathbf{y} = \mathbf{A}\mathbf{x}$ and let $\mathbf{d}_k$ be the $k$-th
generalised diagonal, $d_{k,i} = A_{i,\,i+k \bmod N}$. Then
$y_i = \sum_k d_{k,i} x_{i+k}$, so
$\mathbf{y} = \sum_{k} \mathbf{d}_k \odot \mathrm{rot}_k(\mathbf{x})$, a sum of
$N$ componentwise products of a plaintext with a rotation. For the second bound,
write $k = k_1 \sqrt{N} + k_2$ and factor
$\mathrm{rot}_k = \mathrm{rot}_{k_1\sqrt N} \circ \mathrm{rot}_{k_2}$: precompute
the $\sqrt N$ inner rotations once, and apply $\sqrt N$ outer rotations to the
accumulated partial sums, giving $2\sqrt N$ rotations in total. Correctness is the
associativity of the shift action.
\end{proof}

Proposition~\ref{prop:hs-matvec} is the technique the compilers we compare against
implement and optimise. We record it here because it makes the comparison legible:
when we report that a homomorphic layer costs a fixed number of ciphertext
operations while a secret-sharing layer costs traffic proportional to width, the
homomorphic side's flatness is exactly this packing, and it is also why the
homomorphic side's advantage grows with width in our measurements.

\begin{remark}
\label{rem:ckks}
CKKS encodes \emph{approximate} complex or real vectors rather than exact
integers, using the canonical embedding: a vector is scaled by a factor $\Delta$,
rounded into the ring, and the rounding error is absorbed as additional noise.
Its distinguishing operation is \ib{rescaling}, which divides a ciphertext by
$\Delta$ after a multiplication and is implemented as modulus switching; the
consequence is that CKKS's ``noise'' and its ``precision'' are the same quantity,
and depth is consumed by the scale management rather than tracked separately. Two
properties matter for our comparison. The scheme has no comparison operation and
cannot compute a sign, so a nonlinearity must be a polynomial, which is why the
only activation both paradigms evaluate exactly is a monomial
(Remark~\ref{rem:x2}). And precision degrades along the modulus chain, so a
classification margin narrower than the accumulated error is not recoverable---a
failure mode that presents as a wrong label and not as as an error signal, which
is why our harness compares against a plaintext reference before quoting any
timing.
\end{remark}

\begin{remark}
\label{rem:threshold}
In a \ib{threshold} homomorphic scheme the secret key is shared and decryption is
a protocol. The subtlety is that a partial decryption reveals a noise term that
depends on the key share, so partial decryptions must be flooded by noise drawn
from a range at least $2^{\kappa}$ times the bound---Lemma~\ref{lem:smudge}---and
that flooding is subtracted from the same budget as the evaluation. The practical
consequence, which we report as a disclosed gap and not a solved problem, is
that a statistical security level for the decryption and a depth budget for the
computation are drawn on the same account.
\end{remark}

\subsection{Succinct proofs}
\label{sec:found:zk}

Section~\ref{sec:prelim:zk} describes the system we implement; we give the
definitions here so that the claims in that section are checkable.

\begin{definition}
\label{def:argument}
Let $R$ be a relation and $L = \{x : \exists w,\ (x,w) \in R\}$. An
\ib{interactive argument} for $R$ is a pair $(P,V)$ of interactive machines with
$V$ polynomial-time, satisfying \ib{completeness}---for $(x,w) \in R$, $V$ accepts
with probability $1 - \negl$---and \ib{computational soundness}: for $x \notin L$
and every polynomial-time $P^*$, $V$ accepts with negligible probability. It is
an \ib{argument of knowledge} when there is a polynomial-time extractor that,
given oracle access to a prover convincing $V$ with noticeable probability,
outputs a witness. It is \ib{zero-knowledge} when there is a simulator producing
accepting transcripts indistinguishable from real ones without a witness.
\end{definition}

\begin{definition}
\label{def:snark}
A \ib{SNARK} is a non-interactive argument of knowledge whose proof length and
verification time are $o(\abs{w})$, typically polylogarithmic in the circuit size.
\end{definition}

Two points of principle bear on how we report our zero-knowledge measurements.
First, soundness in Definition~\ref{def:argument} is \emph{computational}: an
unbounded prover can cheat, which is the price of succinctness, and it is why
these objects are called arguments not proofs. Second, non-interactivity
is normally obtained by the Fiat--Shamir transform, replacing the verifier's
random challenges with a hash of the transcript; the transform is provable in the
random-oracle model and its concrete instantiation is an assumption, so a proof
system's security cannot be quoted independently of the hash it is instantiated
with.

\begin{remark}
\label{rem:zk-lesson}
The gap between ``a proof verifies'' and ``a proof proves the intended statement''
is where our own zero-knowledge error lived: a backend can be complete and
internally consistent while the constraint system it emits fails to bind the
witness to the computation one intended, and no amount of successful verification
detects it. The methodological consequence is that a proof system requires a
\ib{negative} test---a witness that ought to fail and does---and that a report of
successful verifications alone is not evidence of soundness.
Appendix~\ref{sec:appendix:caught} records an instance.
\end{remark}
